\documentclass[conference]{IEEEtran}
\IEEEoverridecommandlockouts
\usepackage{cite}
\usepackage{amsmath,amssymb,amsfonts}
\usepackage{graphicx}
\usepackage{textcomp}
\usepackage{url}
\usepackage{caption}
\usepackage{amsmath,amssymb}

\begin{document}

\title{Smart Grid Sensor Network for Industrial Energy\\
Management Using IoT and Machine Learning}

\author{\IEEEauthorblockN{Saeed Sarwar Anas}
\IEEEauthorblockA{\textit{School of Computing, Engineering and Intelligent Systems}\\
\textit{Ulster University}\\
Student Number: 20068400\\
Supervisor: Dr Anwar Haq}}

\maketitle
\begin{abstract}
Industrial facilities account for a substantial share of global electricity
demand, yet most small and medium-sized enterprises (SMEs) operate without any
real-time visibility of their energy use because commercial monitoring
infrastructure is prohibitively expensive. This paper presents the design and
evaluation of a low-cost smart grid sensor network for industrial energy
management that integrates Internet of Things (IoT) sensing with machine
learning. The system couples an inexpensive embedded sensing platform with a
reproducible data pipeline, a Long Short-Term Memory (LSTM) model for
short-term demand forecasting, and an Isolation Forest model for unsupervised
anomaly detection, all surfaced through a real-time dashboard. Using the UCI
Individual Household Electric Power Consumption dataset, the LSTM forecaster
achieved a mean absolute error (MAE) of 0.3707~kW, outperforming a naive
baseline by 8.8\% and a 24-hour moving-average baseline by 34.9\% on unseen
test data. Cross-dataset validation on the Tetouan City load confirmed that the unchanged
pipeline transfers across scales, improving on the same baselines by 62.0\% and
87.0\%. The contribution is an affordable, open-source, end-to-end pipeline
that lowers the barrier to intelligent energy management for resource-
constrained organisations.
\end{abstract}

\begin{IEEEkeywords}
smart grid, Internet of Things, machine learning, LSTM, energy demand
forecasting, anomaly detection, isolation forest, edge computing, Industry~4.0
\end{IEEEkeywords}

\section{Introduction}
The electricity grid is among the most important engineering systems of the
modern era, but for most of the twentieth century it operated as an essentially
``one-way'' network. Electricity flowed from a small number of large, central
generators to passive consumers, and the grid carried almost no information back
about how, when, or where energy was being used. This architecture was robust
and well suited to the demand patterns of its time, yet it left operators
fundamentally blind to consumption at the point of use and unable to respond to
changing conditions in real time.

Several converging pressures over the past two decades have exposed the limits
of this traditional model. Global electricity demand has continued to rise, with
industry responsible for a large fraction of total consumption
\cite{international2022international}. At the same time, decarbonisation targets
and the rapid uptake of intermittent renewable generation have demanded a far
more flexible and observable grid. These pressures gave rise to the concept of
the \emph{smart grid}: an electricity network augmented with digital sensing,
two-way communication, and data-driven control, so that supply and demand can be
monitored and balanced dynamically rather than managed reactively. Surveys of
smart grid communication infrastructure document this shift from a passive
network towards a sensor-rich, communication-enabled system
\cite{fan2012smart,yan2012survey}.

Three paradigms have shaped the sector since. The first is the move from
centralised generation towards \emph{distributed} energy resources, which
multiplies the points that must be observed and coordinated
\cite{silva2023iot}. The second is \emph{Industry~4.0} and the Industrial
Internet of Things, in which low-cost networked sensors make pervasive,
fine-grained measurement economically feasible
\cite{munirathinam2020industry}. The third is \emph{machine learning}, which
has moved energy systems from threshold-based monitoring towards predictive
management: deep learning models non-linear patterns in sequential data
\cite{lecun2015deep}, recurrent architectures such as the Long Short-Term
Memory (LSTM) network \cite{hochreiter1997long} match the daily and weekly
cycles of consumption, and unsupervised methods such as the Isolation Forest
\cite{liu2008isolation} flag abnormal usage without the labelled fault data
that is rarely available in practice. The combined impact is significant:
sensor-driven feedback has been shown to reduce consumption
\cite{alaa2017review,hafeez2020efficient}, LSTM forecasters to outperform
classical statistical baselines \cite{kong2017short,raza2015review}, and
anomaly detection to surface equipment faults before they escalate
\cite{himeur2021artificial}.

Despite this progress, a clear gap remains. Each of the four integrated
energy-management platforms reviewed in Section~\ref{sec:litreview}
\cite{silva2023iot,vijayalakshmi2023smart,bin2022real,ahmed2020machine}
depends on proprietary or edge--cloud infrastructure---supervisory control and
data acquisition (SCADA) systems,
industrial programmable logic controllers, and commercial smart meters whose
cost runs to many thousands of pounds. This places intelligent energy management
firmly out of reach of small and medium-sized enterprises (SMEs), even though
SMEs collectively account for a large share of industrial energy use. Moreover,
much of the existing research addresses individual components in isolation a
forecasting model here \cite{wen2024deep}, an anomaly detector there
\cite{pereira2018unsupervised} rather than demonstrating a complete,
reproducible pipeline. As a result, an affordable and openly reproducible
end-to-end system that an under-resourced organisation could actually adopt has
remained largely unaddressed.

This paper addresses that gap. It presents the design, implementation, and
evaluation of a low-cost smart grid sensor network for industrial energy
management that integrates IoT sensing with machine learning in a single,
reproducible pipeline. In contrast to prior work that depends on costly
proprietary hardware or that treats forecasting and anomaly detection
separately, the proposed system combines inexpensive embedded sensing (targeting
hardware of the order of \pounds30), a documented and version-controlled data
pipeline, an LSTM-based demand forecaster, an Isolation~Forest anomaly detector,
and a real-time monitoring dashboard. The work is grounded in a Design Science
Research methodology \cite{hevner2004design,peffers2007design} and follows
established guidance for the conduct of the research project
\cite{zina2021essential}.

\noindent The principal contributions of this paper are:
\begin{itemize}
  \item an affordable, open-source, end-to-end architecture that brings together
        low-cost IoT sensing, wireless telemetry, time-series forecasting,
        unsupervised anomaly detection, and real-time visualisation;
  \item an empirical evaluation of an LSTM demand forecaster on a public
        benchmark dataset, including a comparison against naive and
        moving-average baselines; and
  \item a reproducible pipeline---with code, processed data, and trained models
        released publicly that lowers the barrier to intelligent energy
        management for SMEs.
\end{itemize}

The remainder of this paper is organised as follows. Section~\ref{sec:problem}
defines the problem and its scale. Section~\ref{sec:aims} states the aim and
objectives. Section~\ref{sec:litreview} reviews related work, critically
compares existing approaches, and derives the set of parameters that informs the
proposed solution architecture. Section~\ref{sec:method} describes the
methodology and system design, Section~\ref{sec:results} reports and discusses
the results, and Section~\ref{sec:conclusion} concludes and outlines future
work.

\section{Problem Statement}\label{sec:problem}
Industrial activity accounts for approximately one third of total global energy
use \cite{international2022international}. A substantial part is consumed
inefficiently, and in many facilities that inefficiency is never detected
because consumption is not measured at the level of individual processes or
machines.
 
The burden falls disproportionately on SMEs, which constitute the majority of
industrial and commercial organisations yet mostly operate without intelligent
monitoring. The principal barrier is cost: SCADA systems, industrial
programmable logic controllers, and commercial smart-metering platforms
typically require capital investment running to tens of thousands of pounds
alongside specialist installation and maintenance
\cite{vijayalakshmi2023smart,silva2023iot}. The consequence is that these
businesses operate effectively ``blind'': waste accumulates unseen, faults
develop without early warning, and there is no basis on which to forecast
demand.
 
Where monitoring exists it is frequently reactive. Threshold-based alarms
register a problem only after a value has crossed a fixed limit and are poorly
suited to the gradual drift or context-dependent anomalies that precede failure
\cite{himeur2021artificial}. Advances in IoT sensing and machine-learning
forecasting indicate that a proactive approach is achievable, but have largely
been demonstrated as isolated components rather than a deployable system.
 
The core problem is therefore the absence of an affordable, integrated, and reproducible system that can sense conditions in real time, forecast short-term demand, and detect anomalous consumption within a budget accessible to an SME;
the platforms that come closest to this integration achieve it only through edge--cloud infrastructure whose cost reintroduces precisely the barrier in question \cite{bin2022real,ahmed2020machine}.
It is non-trivial because a credible solution must integrate embedded sensing,
time-series deep learning, unsupervised detection, and reproducible data
engineering into one coherent system, then be evaluated rigorously against
established baselines rather than asserted to work.
 
\section{Aim and Objectives}\label{sec:aims}
 
\textit{Aim.} To design, implement, and evaluate a low-cost smart grid sensor
network for industrial energy management that uses the Internet of Things and
machine learning to deliver real-time monitoring, short-term demand
forecasting, and anomaly detection, thereby demonstrating that capabilities
previously confined to expensive proprietary systems can be achieved with
affordable hardware and open-source software.
 
\noindent The objectives that operationalise this aim are as follows:
 
\begin{enumerate}
\item To critically review the literature on IoT-based energy monitoring,
machine-learning demand forecasting, and anomaly detection, in order to identify
the research gap and define a precise project scope.
 
\item To acquire and prepare benchmark energy data principally the UCI
Individual Household Electric Power Consumption dataset
\cite{hebrail2012individual} through a documented, version-controlled
preprocessing pipeline implemented in Python with \texttt{pandas}, applying
chronological train/validation/test partitioning to prevent data leakage.
 
\item To design the IoT sensing architecture around a low-cost embedded platform, using Bluetooth Low
Energy and Wi-Fi telemetry to stream environmental measurements to a time-series
datastore \cite{munirathinam2020industry,warden2019tinyml}.
 
\item To develop and train a Long Short-Term Memory (LSTM) network
\cite{hochreiter1997long} for short-term energy demand forecasting using
TensorFlow/Keras \cite{abadi2016tensorflow}, and to benchmark it against naive
and moving-average baselines \cite{makridakis2018statistical}.
 
\item To implement unsupervised anomaly detection using an Isolation Forest model
\cite{liu2008isolation} that flags abnormal consumption patterns without
requiring labelled fault data.

\item To build a real-time monitoring dashboard using Streamlit that visualises
live readings, forecasted demand, and anomaly alerts in a form usable by a
non-specialist operator.
 
\item To evaluate all models rigorously on held-out test data using appropriate
quantitative metrics---mean absolute error (MAE), root-mean-square error (RMSE),
and mean absolute percentage error (MAPE) for forecasting, and precision,
recall, and F1-score for anomaly detection.
 
 
\end{enumerate}

\section{Literature Review}\label{sec:litreview}
This work draws on four overlapping bodies of research: IoT-based energy
monitoring, machine-learning demand forecasting, anomaly detection, and
low-cost embedded computing.
 
\subsection{IoT-Based Energy Monitoring}
Broad surveys of IoT-enabled smart grids
\cite{goudarzi2022survey} establish that sensing,
communication, and analytics are now treated as a single integrated stack. At
the application level, reviews of smart-home IoT \cite{alaa2017review} and
smart-meter analytics \cite{yildiz2017recent} show that fine-grained
measurement with feedback influences consumption, while Hafeez \textit{et al.}
\cite{hafeez2020efficient} couple sensing with price-based demand response.
At system scale, Silva
\textit{et al.} \cite{silva2023iot} address AC microgrids under grid and
security constraints, Vijayalakshmi \textit{et al.}
\cite{vijayalakshmi2023smart} demonstrate energy management on embedded
systems, and Bin Mofidul \textit{et al.} \cite{bin2022real} combine
acquisition, anomaly detection, and monitoring across edge and cloud AI. Two limitations recur: attention concentrates on individual residences and
utility-scale systems, leaving the industrial SME underserved; and the more
capable integrated platforms achieve breadth through costly edge--cloud
infrastructure, reintroducing the affordability barrier.
 
\subsection{Machine Learning for Demand Forecasting}
The field has shifted decisively from statistical models towards deep learning
\cite{lecun2015deep}, with reviews tracing
rapid adoption for load forecasting \cite{raza2015review}.
The LSTM \cite{hochreiter1997long} has become dominant for short-term load:
Kong \textit{et al.} \cite{kong2017short} demonstrated its ability to capture
household dynamics and Inala \textit{et al.} \cite{inala2025enhancing} apply
the same family in a smart-grid setting. Comparative studies refine this
picture---Mahjoub \textit{et al.} \cite{mahjoub2022predicting} benchmark LSTM
against multi-layer and drop-GRU variants and expose an accuracy--efficiency
trade-off, while Wen \textit{et al.} \cite{wen2024deep} pursue accuracy through
a hybrid model at the cost of complexity. Crucially, Makridakis \textit{et al.}
\cite{makridakis2018statistical} found that machine-learning forecasters do not
automatically outperform simpler methods and argued for evaluation against
strong baselines---a finding treated here as binding rather than advisory. The
broader gap is that most studies optimise accuracy on a single dataset and
report little on deployability: model size, inference cost, or suitability for
constrained hardware.
 
\subsection{Anomaly Detection in Energy Systems}
Himeur \textit{et al.} \cite{himeur2021artificial} review AI-based anomaly
detection in buildings and identify label scarcity as the field's central
tension, dividing the literature into two camps. Supervised approaches---
XGBoost with LSTM \cite{chen2024effective}, comparative classifier studies, and network-level detection \cite{sahu2020machine}---
can be highly accurate but presuppose labelled examples of the very faults they
aim to catch. Unsupervised methods avoid that dependency: the Isolation Forest
\cite{liu2008isolation} isolates anomalies efficiently, clustering has been
applied to energy series \cite{jesmeen2021unsupervised}, and variational
recurrent autoencoders achieve strong results at considerable computational
weight \cite{pereira2018unsupervised}. Decisively for this work, Jamshidi
\textit{et al.} \cite{jamshidi2025lightweight} pair an autoencoder with an
Isolation Forest specifically for resource-constrained IoT edge gateways,
confirming the Isolation Forest belongs at the lightweight end of the spectrum.
 
\subsection{Low-Cost Embedded and Edge Computing}
TinyML \cite{warden2019tinyml} established that machine-learning inference can
run on low-power microcontrollers, and Applied studies confirm this:
most directly relevant, Baskar \textit{et al.} \cite{baskar2024edge} deploy the Arduino Nano 33 BLE Sense, the platform
adopted here, for industrial wastewater monitoring. The remaining gap is
integrative: edge-ML studies typically address a single task and rarely combine
forecasting, detection, and visualisation into one reproducible system.
 
\subsection{Synthesis and Proposed Architecture}
Taken together, the literature establishes four points. IoT monitoring is
mature but concentrated away from the industrial SME. The LSTM is effective but
must be validated against simple baselines and is rarely assessed for
deployability. Lightweight unsupervised detection is the natural fit where
labels are scarce and compute is limited. Low-cost edge hardware is proven for
sensing. What remains absent is a single affordable, openly reproducible system
integrating all four; Ahmed \textit{et al.} \cite{ahmed2020machine} move toward
such integration but not at the SME-accessible price point this project
targets.
 
These findings yield the design parameters that govern the system:
affordability, realised as a $\sim$\pounds30 embedded platform; baseline
validation, realised as mandatory naive and moving-average comparison; a
label-free anomaly paradigm, realised as an Isolation Forest chosen over deep
autoencoders for its computational footprint; leakage-free evaluation, realised
as a chronological split with training-only scaler fitting; and
reproducibility, realised as public release of code, data, and models. Guided
by these, the proposed system is the four-layer architecture of
Fig.~\ref{fig:arch}: an IoT sensing layer on the validated low-cost platform
\cite{baskar2024edge}; a data-pipeline layer providing leakage-free
preprocessing; a machine-learning layer pairing an LSTM forecaster
\cite{kong2017short} with an Isolation Forest detector \cite{liu2008isolation};
and a visualisation layer delivering a real-time dashboard. Compared with the
edge--cloud system of Bin Mofidul \textit{et al.} \cite{bin2022real}, the
design deliberately trades infrastructural sophistication for affordability and reproducibility.
\begin{figure}[t]
  \centering
  \includegraphics[width=\columnwidth]{images/Architecture_Professor_Format.png}
  \caption{Proposed four-layer system architecture.}
  \label{fig:arch}
\end{figure}
\section{Methodology and System Design}\label{sec:method}
\subsection{Research Approach}
The project follows a Design Science Research (DSR) methodology
\cite{hevner2004design,peffers2007design}, in which knowledge is generated
through the construction and rigorous evaluation of an artefact here, the
integrated smart grid monitoring system. Because the requirements were
exploratory and the work was intense in research, an incremental 
and research process was adopted, comprising four phases: data acquisition and
preparation, model development, evaluation against baselines and integration of
the models into a monitoring dashboard. The physical sensing layer follows the
low-cost architecture set out in Section~\ref{sec:litreview}; the evaluation
reported here uses a public benchmark dataset as its data source, with live
hardware deployment reserved for future work.
 
\subsection{Dataset and Preprocessing}
The secondary dataset is the UCI Individual Household Electric Power Consumption
dataset \cite{hebrail2012individual}, comprising 2{,}075{,}259 minute-level
measurements recorded between December 2006 and November 2010. It contains seven
electrical quantities: global active power, global reactive power, voltage,
global intensity, and three sub-metering channels. This dataset was selected
because it provides a long, continuous, real-world record whose daily, weekly,
and seasonal structure is well suited to sequence modelling. It is worth noting
that other public benchmarks target related but distinct problems; the Smart
Grid Stability dataset \cite{arzamasov2018towards}, for example, poses a binary
stability-classification task rather than continuous demand forecasting, and is
therefore outside the scope of the forecasting model developed here.
 
The raw data were prepared through a documented, version-controlled pipeline.
First, 25{,}979 missing values were imputed by forward-fill, the standard
choice for time series in which the most recent observation best estimates the
next. Second, the series was resampled from one-minute to hourly resolution,
producing 34{,}589 clean records and reducing noise while preserving the
temporal patterns relevant to energy management. Third, extreme values were
screened using the interquartile-range method, and all features were normalised
to $[0,1]$ using min--max scaling as required for stable recurrent-network
training. Fourth, a sliding window of 24 hours was applied, yielding 34{,}565
supervised input--output sequences. Finally, the data were partitioned
chronologically into training, validation, and test sets in a 70/15/15 ratio
(24{,}195, 5{,}184, and 5{,}186 samples respectively), with the scaler fitted on
the training partition only so that no information from future time steps could
leak into the model input \cite{makridakis2018statistical}.
 
\begin{equation}
x' = \frac{x - x_{\min}}{x_{\max} - x_{\min}}, \label{eq:minmax}
\end{equation}
 
where $x_{\min}$ and $x_{\max}$ are computed on the training partition only. The
forecasting task was then framed as a supervised sliding-window problem: given
the previous $w = 24$ hourly observations, the model predicts the next hour,
 
\begin{equation}
\hat{y}_{t} = f\!\left(x_{t-w}, x_{t-w+1}, \dots, x_{t-1}\right), \label{eq:window}
\end{equation}
 
where $f(\cdot)$ is the trained LSTM.
 
\subsection{Demand Forecasting Model}
Short-term demand forecasting was addressed with a Long Short-Term Memory (LSTM)
network \cite{hochreiter1997long}, whose gating mechanism is well matched to the
multi-scale temporal dependencies in energy data and which has an established
record in load forecasting \cite{kong2017short}. At each time step $t$, an LSTM unit updates its cell state $C_t$ and hidden
state $h_t$ from the current input $x_t$ and the previous hidden state
$h_{t-1}$ through three gates---forget ($f_t$), input ($i_t$), and output
($o_t$)---defined as
 
\begin{align}
f_t &= \sigma\!\left(W_f \cdot [h_{t-1}, x_t] + b_f\right), \label{eq:forget}\\
i_t &= \sigma\!\left(W_i \cdot [h_{t-1}, x_t] + b_i\right), \label{eq:input}\\
\tilde{C}_t &= \tanh\!\left(W_C \cdot [h_{t-1}, x_t] + b_C\right), \label{eq:cand}\\
C_t &= f_t \odot C_{t-1} + i_t \odot \tilde{C}_t, \label{eq:cell}\\
o_t &= \sigma\!\left(W_o \cdot [h_{t-1}, x_t] + b_o\right), \label{eq:output}\\
h_t &= o_t \odot \tanh(C_t), \label{eq:hidden}
\end{align}
 
where $W_{\{f,i,C,o\}}$ and $b_{\{f,i,C,o\}}$ are the learned weight matrices and
bias vectors, $\sigma(\cdot)$ is the logistic sigmoid, $\tanh(\cdot)$ is the
hyperbolic tangent, and $\odot$ denotes element-wise multiplication. The forget
gate \eqref{eq:forget} controls how much of the previous cell state is retained,
the input gate \eqref{eq:input} and candidate \eqref{eq:cand} determine what new
information is written, and the output gate \eqref{eq:output} produces the hidden
state passed to the next step. This gating is what allows the network to retain
information over the long daily and weekly cycles present in the energy series. The network comprises two
stacked LSTM layers of 64 and 32 units, each followed by a dropout layer with
rate 0.2, and a dense output neuron, giving 29{,}857 trainable parameters. It
was trained to minimise the mean-squared-error loss using the Adam optimiser
with a batch size of 32 for up to 50 epochs, with early stopping (patience of 10
epochs) to prevent overfitting; the implementation used TensorFlow/Keras
\cite{abadi2016tensorflow} with a fixed random seed for reproducibility. To place
the model's performance in context, two transparent baselines were implemented:
a naive persistence predictor, which forecasts the next hour as equal to the
current one, and a 24-hour moving average.
 
\subsection{Anomaly Detection Model}
Anomalous consumption was identified with an Isolation Forest
\cite{liu2008isolation}, chosen because it is unsupervised---and therefore
requires no labelled fault data, which is unavailable for this dataset---and
because it is computationally lightweight and thus suited to low-cost
deployment. The method builds an ensemble of trees that recursively partition
the feature space along randomly selected features and split points, scoring
each observation by the mean path length required to isolate it: anomalies,
being few and readily separable, isolate in short paths, and a threshold set by
the contamination parameter separates them from the remainder.

The detector operated on all seven electrical features simultaneously (a
multivariate approach), after standardising each feature to zero mean and unit
variance so that no single quantity dominated by virtue of its scale. The model
was configured with 200 base estimators, a contamination rate of 0.01,
automatic subsampling, and a fixed random seed. Because the dataset carries no
ground-truth anomaly labels, a transparent statistical baseline---a Z-score
rule flagging any hour whose global active power lies more than three standard
deviations from the mean---was used as a reference point for comparison.
 
\subsection{Evaluation Metrics}
Forecasting accuracy was quantified on the held-out test set using the mean
absolute error (MAE) and root-mean-square error (RMSE), both expressed in
kilowatts for direct interpretability, together with the mean absolute
percentage error (MAPE) and the coefficient of determination $R^2$; the LSTM
was compared against the naive and moving-average baselines on identical data.
For anomaly detection, in the absence of ground-truth labels, precision,
recall, and the F1-score were computed \emph{relative to the Z-score
reference}, and therefore measure agreement between the two methods rather than
correctness against true anomalies. They are complemented by an analysis of the
multivariate anomalies the Isolation Forest detects but the univariate baseline
misses.
 
\begin{align}
\mathrm{MAE} &= \frac{1}{n}\sum_{t=1}^{n}\left|y_t - \hat{y}_t\right|, \label{eq:mae}\\
\mathrm{RMSE} &= \sqrt{\frac{1}{n}\sum_{t=1}^{n}\left(y_t - \hat{y}_t\right)^2}, \label{eq:rmse}\\
\mathrm{MAPE} &= \frac{100}{n}\sum_{t=1}^{n}\left|\frac{y_t - \hat{y}_t}{y_t}\right|, \label{eq:mape}\\
R^2 &= 1 - \frac{\sum_{t=1}^{n}(y_t-\hat{y}_t)^2}{\sum_{t=1}^{n}(y_t-\bar{y})^2}, \label{eq:r2}
\end{align}
 
where $y_t$ is the actual value, $\hat{y}_t$ the prediction, $\bar{y}$ the mean
of the actual values, and $n$ the number of test observations. For the anomaly
detector, agreement with the Z-score reference was quantified using precision,
recall, and their harmonic mean, the F1-score,
 
\begin{equation}
P = \frac{TP}{TP+FP}, \quad
R = \frac{TP}{TP+FN}, \quad
F_1 = \frac{2PR}{P+R}, \label{eq:prf}
\end{equation}
 
where $TP$, $FP$, and $FN$ are the counts of records flagged by both methods, by
the Isolation Forest only, and by the Z-score reference only, respectively. The
Z-score used by the baseline is
 
\begin{equation}
z_t = \frac{x_t - \mu}{\sigma}, \label{eq:zscore}
\end{equation}
 
with $\mu$ and $\sigma$ the mean and standard deviation of global active power;
an hour is flagged when $|z_t| > 3$.
 
\subsection{Implementation and Reproducibility}
The system was implemented in Python~3.11 using TensorFlow/Keras, scikit-learn,
pandas, and NumPy. Development and training were carried out on a consumer laptop
(AMD Ryzen~7 6800U, 16~GB RAM), with the option of GPU-accelerated training on a
cloud platform for larger workloads. The visualisation layer is a Streamlit
dashboard that loads the saved model outputs and presents demand forecasts
alongside anomaly alerts. To support full reproducibility, fixed random seeds
were used throughout, and all source code, the processed dataset, trained model
files, and results have been released in a public repository.

\section{Results and Discussion}\label{sec:results}
This section reports the performance of the two models on the held-out data and
then discusses their practical significance and limitations. Unless otherwise
stated, forecasting metrics are reported over a 500-hour window of the held-out
test set, consistent with the figures presented in the monitoring dashboard.
 
\subsection{Demand Forecasting Performance}
Table~\ref{tab:forecast} reports the LSTM alongside the two baselines. The
LSTM achieved the lowest error on the metrics appropriate to this task, with an
MAE of 0.3707~kW and an RMSE of 0.5421~kW---within roughly 370~W of true hourly
demand---reducing MAE by 8.6\% against naive persistence and 35.7\% against the
moving average, and explaining 44.3\% of the variance ($R^2 = 0.4432$) against
23.6\% and 5.4\% respectively. Figure~\ref{fig:uciforecast} shows the forecast
tracking the actual series.  
 
\begin{table}[t]
\centering
\caption{Forecasting performance on the held-out test set (lower is better, except $R^2$)}
\label{tab:forecast}

\resizebox{\columnwidth}{!}{%
\begin{tabular}{|l|c|c|c|c|}
\hline
\textbf{Model} & \textbf{MAE (kW)} & \textbf{RMSE (kW)} & \textbf{MAPE (\%)} & \textbf{$R^2$} \\
\hline
Naive & 0.4054 & 0.6350 & 44.07 & 0.2361 \\
\hline
Moving Average & 0.5768 & 0.7066 & 87.89 & 0.0542 \\
\hline
\textbf{LSTM} & \textbf{0.3707} & \textbf{0.5421} & 45.77 & \textbf{0.4432} \\
\hline
\end{tabular}
}
\end{table}
\begin{figure}[t]
  \centering
  \includegraphics[width=\columnwidth]{images/04_actual_vs_predicted.png}
  \caption{LSTM forecast against actual demand on the UCI test set.}
  \label{fig:uciforecast}
\end{figure}
 
The LSTM's MAPE of 45.77\% marginally exceeds naive's 44.07\% despite clear
superiority on every other measure. This reflects a known weakness of MAPE on
such data: household consumption falls close to zero overnight, and dividing by
near-zero values inflates percentage error irrespective of absolute accuracy.
MAPE is therefore not a reliable discriminator here. The moving average
performs worst throughout, its smoothing lagging the peaks the LSTM
anticipates.
 
\subsection{Anomaly Detection Performance}
The Isolation Forest, operating on all seven electrical features, flagged 346
of the 34{,}589 hourly records (1.00\%) as anomalous;
Figure~\ref{fig:ucianom} shows their distribution across the series and in
score space. Against the univariate Z-score baseline, which flagged 440
records, the methods agreed on 267, the Isolation Forest alone flagged 79, and
the Z-score rule alone flagged 173 (Table~\ref{tab:anomaly}).
 
\begin{table}[t]
\centering
\caption{Agreement between the Isolation Forest and the Z-score baseline}
\label{tab:anomaly}
\small
\renewcommand{\arraystretch}{1.3}
\begin{tabular}{|l|c|}
\hline
\textbf{Category} & \textbf{Records}\\
\hline
Flagged by Isolation Forest & 346\\
\hline
Flagged by Z-score baseline & 440\\
\hline
Flagged by both & 267\\
\hline
Isolation Forest only (multivariate) & 79\\
\hline
Z-score only (univariate) & 173\\
\hline
\end{tabular}
\end{table}
\begin{figure}[t]
  \centering
  \includegraphics[width=\columnwidth]{images/05_anomaly_detection.png}
  \caption{Detected anomalies across the UCI series (top) and the anomaly-score distribution (bottom).}
  \label{fig:ucianom}
\end{figure}
 
Treating the Z-score flags as a reference gives a precision of 0.7717, a recall
of 0.6068, and an F1-score of 0.6794. These figures measure \emph{agreement}
between two methods rather than correctness against true anomalies, since no
ground-truth labels exist; the Z-score rule is a transparent reference, not an
authority. Read this way, the Isolation Forest reproduces the anomalies a
simple statistical rule catches, while its distinctive value lies in the 79
additional records---unusual \emph{combinations} such as elevated intensity and
sub-metering activity without an extreme total-power value, which a
single-feature rule cannot detect by construction.

\subsection{Cross-Dataset Validation on the Tetouan City Dataset}
To assess generalisation beyond a single household, the complete workflow was
re-applied, unchanged, to the Power Consumption of Tetouan City dataset:
52{,}416 ten-minute records for 2017 across three distribution zones of
Tetouan, Morocco, with meteorological variables. The data were resampled to
hourly resolution (8{,}736 records, no missing values), zone loads aggregated
into total city consumption, and the same 24-hour window and chronological
70/15/15 split applied (6{,}115 training, 1{,}310 validation, 1{,}311 test
samples) with identical model configuration.

Table~\ref{tab:tetouan} reports the results. The LSTM achieved an MAE of
1{,}521.60~kW against 4{,}007.33~kW for naive persistence and 11{,}695.04~kW
for the moving average---improvements of 62.0\% and 87.0\%---with a MAPE of
2.41\% and a normalised MAE of 2.36\% of mean load;
Figure~\ref{fig:tetforecast} shows the close tracking of the smooth city-scale
series. Accuracy is markedly higher than on household data, as expected: a
city-scale load aggregates thousands of consumers into a smooth, strongly
periodic series, whereas a single household is dominated by irregular appliance
events. MAPE is also reliable here because city-scale consumption never
approaches zero.

\begin{table}[t]
\centering
\caption{Forecasting performance on the Tetouan City test set}
\label{tab:tetouan}
\small
\resizebox{\columnwidth}{!}{%
\begin{tabular}{|l|c|c|c|c|}
\hline
\textbf{Model} & \textbf{MAE (kW)} & \textbf{RMSE (kW)} &
\textbf{MAPE (\%)} & \textbf{nMAE (\%)}\\
\hline
Naive & 4{,}007.33 & 5{,}465.53 & 6.24 & 6.21\\
\hline
Moving Average & 11{,}695.04 & 14{,}034.60 & 18.92 & 18.12\\
\hline
\textbf{LSTM} & \textbf{1{,}521.60} & \textbf{2{,}299.07} & \textbf{2.41} &
\textbf{2.36}\\
\hline
\end{tabular}
}
\end{table}

\begin{figure}[t]
  \centering
  \includegraphics[width=\columnwidth]{images/06_tetouan_forecast.png}
  \caption{LSTM forecast against actual city-scale demand on the Tetouan test
  set.}
  \label{fig:tetforecast}
\end{figure}

The detector was likewise re-applied, operating multivariately on total
consumption with the meteorological features. It flagged 88 of the 8{,}736
records, of which only 7 coincided with the 35 flagged by the Z-score rule
(F1 $= 0.1138$). This low agreement is expected by construction rather than
indicative of failure: a smooth seasonal load produces few $|z|>3$ extremes,
while the multivariate detector responds to unusual \emph{weather--demand
combinations} invisible to a consumption-only threshold. Consistent with this,
the 81 records flagged only by the Isolation Forest concentrate in July and
August during afternoon and evening hours---heat-driven peaks in which high
temperature and high load co-occur atypically, visible in
Figure~\ref{fig:tetanom}.
\begin{figure}[t]
  \centering
  \includegraphics[width=\columnwidth]{images/07_tetouan_anomalies.png}
  \caption{Isolation Forest detections on the Tetouan series, concentrated in
  summer demand peaks.}
  \label{fig:tetanom}
\end{figure} 
\subsection{Discussion}
The results support the central claim of this work: that capable energy
analytics can be delivered with lightweight, low-cost methods rather than
proprietary infrastructure. Three findings deserve emphasis.

First, the improvement over baselines is modest on household data and
substantial on aggregated data, and this asymmetry is informative rather than
incidental. A single household is dominated by stochastic appliance events that
no model can anticipate from load history alone, so the achievable ceiling is
low; a city-scale series aggregates thousands of such events into a smooth,
strongly periodic signal in which the LSTM's capacity to model multi-scale
temporal structure is fully exercised. The practical implication for an SME is
encouraging: an industrial facility with many machines running to a production
schedule sits closer to the aggregated regime than to the single-household one,
so the accuracy demonstrated on Tetouan is the more representative indication
of what such a deployment would achieve. This behaviour is consistent with the
established record of the LSTM on short-term load forecasting
\cite{kong2017short} and with the caution of Makridakis \textit{et al.}
\cite{makridakis2018statistical} that superiority must be demonstrated rather
than assumed---a caution treated here as binding, which is why the baselines
were evaluated on identical data rather than quoted from the literature.

Second, the two detection regimes are complementary rather than competing. On
volatile household data the Isolation Forest largely reproduces what a
univariate statistical rule already catches, adding a margin of 79 combination
anomalies; on smooth aggregated data almost its entire contribution is
multivariate, since few observations are extreme in any single feature. A
practitioner should therefore expect the value of multivariate detection to
increase precisely where simple thresholds become least informative---the
regime an industrial deployment is most likely to occupy. That the detector
required no labelled faults is what makes it deployable at all, given that such
labels are rarely available in practice \cite{himeur2021artificial}.

Third, affordability was achieved without a corresponding sacrifice in
capability. The forecaster carries 29{,}857 parameters and trained in
approximately eleven minutes on a consumer laptop CPU with no GPU at any stage;
both models are small enough to be compatible with the low-cost deployment the
project targets, and the complete pipeline runs in a working dashboard.

Several limitations qualify these conclusions. Forecasting metrics are reported
over a bounded test window rather than the full multi-year series, so they
characterise typical rather than worst-case performance. The anomaly
evaluation, lacking labelled faults, is an agreement study rather than a
measurement of detection accuracy, and its precision and recall should not be
read as such. Neither dataset is industrial: the cross-dataset validation
establishes transferability across load scales but not across load
\emph{types}, and validation on live industrial sensor data remains future
work. Finally, both models are black boxes, and post-hoc explanation would be
necessary before an operator could be expected to act on a flagged hour with
confidence. These limitations are revisited in Section~\ref{sec:conclusion}.
\section{Conclusion and Future Work}\label{sec:conclusion}
This paper set out to demonstrate that intelligent energy
management---real-time monitoring, short-term demand forecasting, and anomaly
detection---can be delivered at a cost accessible to small and medium-sized
enterprises, using low-cost hardware and open-source software rather than
proprietary infrastructure. An end-to-end pipeline was designed, implemented,
and evaluated: a documented, leakage-free preparation stage reducing
2{,}075{,}259 raw observations to 34{,}589 hourly records; a two-layer LSTM
forecaster of 29{,}857 parameters benchmarked against transparent baselines; a
multivariate Isolation Forest requiring no labelled faults; and a dashboard
surfacing both. The complete system was released publicly for independent
replication.

The empirical results support the central claim. On the UCI benchmark the LSTM
achieved an MAE of 0.3707~kW---within roughly 370~W of true hourly
demand---reducing error by 8.6\% against naive persistence and 35.7\% against a
24-hour moving average, while explaining 44.3\% of the variance against 23.6\%
for naive. The Isolation Forest flagged 346 records, 267 agreeing with a
Z-score reference and 79 identified by the multivariate detector alone.
Cross-dataset validation on the Tetouan City load confirmed that the unchanged
pipeline transfers across scales, with improvements of 62.0\% and 87.0\% and a
normalised MAE of 2.36\%. The entire evaluation ran on a consumer laptop CPU,
with training completing in approximately eleven minutes---itself part of the
affordability argument.

The work has limitations that define its future direction. The evaluation is
grounded in secondary data from two public benchmarks, which is the scope of
this project; extending it to primary data from the designed sensing layer on
the Arduino Nano 33 BLE Sense platform is a natural direction for further work.
The anomaly evaluation remains an agreement study absent labelled faults, and
deployment feedback would allow the contamination parameter to be tuned against
ground truth. Both models are presently black boxes; SHAP-based attribution
\cite{lundberg2017unified} would let operators see \emph{why} an hour was
flagged. Finally, the forecaster could be extended from univariate to
multivariate input, exploiting the meteorological variables the Tetouan
analysis shows interact strongly with demand.

In sum, the contribution of this work is not a new algorithm but a
demonstration: that established, lightweight methods, engineered carefully into
a single reproducible system, are sufficient to bring smart-grid capabilities
within reach of the organisations that currently lack them.

\bibliographystyle{IEEEtran}
\bibliography{References}

\end{document}